\section{Ограничения и направления развития}

\subsection{Ограничения текущей версии системы}

Несмотря на полученные результаты, разработанную систему следует рассматривать как исследовательский прототип с
ограниченной областью применимости. Эти ограничения связаны как с природой исходных данных, так и с выбранной
постановкой задачи, архитектурой решения и способом интерпретации результатов. Их явное указание важно для корректной
оценки вклада работы и для понимания того, в каких границах полученные выводы остаются обоснованными.

\subsubsection{Ограниченность постановки задачи}

Прежде всего, система не является полноценным оптимизатором севооборота. В работе решается более узкая задача: по
истории культур на конкретном поле и базовым признакам формируется ранжированный список вероятных следующих культур.
Такой подход полезен как инструмент поддержки принятия решений, однако он не охватывает весь набор факторов, влияющих на
выбор культуры в реальном хозяйстве. В частности, в текущей постановке не учитываются:

\begin{itemize}
\tightlist
\item
  экономические показатели;
\item
  погодные сценарии и сезонные прогнозы;
\item
  почвенные ограничения;
\item
  логистические и ресурсные ограничения хозяйства;
\item
  многолетнее планирование на уровне севооборотной схемы.
\end{itemize}

Следовательно, даже качественный shortlist не должен интерпретироваться как окончательное или гарантированно оптимальное
решение. Корректнее рассматривать его как набор вероятных и содержательно интерпретируемых вариантов, которые могут быть
использованы специалистом как исходная база для дальнейшего анализа.

\subsubsection{Ограничения, связанные с данными}

Существенное ограничение связано с используемым пространством целевых классов. После агрегации и фильтрации редких
культур модель обучается и оценивается на наборе из девяти классов. Такое решение было необходимо для стабилизации
постановки и для уменьшения дисбаланса классов, однако оно сокращает полноту исходного пространства культур. В
результате система не охватывает редкие и специализированные случаи, присутствующие в исходных данных.

Кроме того, модель опирается на исторические закономерности, наблюдаемые в имеющихся последовательностях культур. Это
означает, что качество рекомендаций зависит от структуры и покрытия исходного датасета. Если для определённых типов
полей, регионов или культур исторических наблюдений недостаточно, то и модельный результат для них будет менее
устойчивым. В этом смысле система лучше работает в пространстве часто встречающихся и хорошо представленных паттернов,
чем в редких или нетипичных сценариях.

\subsubsection{Ограничения предиктивного ядра}

Хотя baseline CatBoost показывает сильный результат в Top-k постановке, сама модель остаётся статистическим
инструментом, извлекающим зависимости из данных. Она не обладает причинным пониманием агрономической полезности культуры
и не может интерпретировать свои рекомендации как агрономические правила. Это означает, что высокий модельный score
следует понимать как отражение наблюдаемой закономерности в данных, а не как доказательство того, что данная культура
объективно является наилучшей с агрономической точки зрения.

Кроме того, качество модели неодинаково по классам. Для части культур top-1 качество высоко уже в базовой постановке,
тогда как для более сложных классов полезность системы проявляется главным образом в формате shortlist. Это ограничивает
интерпретацию результата: в некоторых случаях система может быть надёжна уже как top-1 предиктор, но в других корректнее
использовать её именно как рекомендательную, а не классификационную систему.

\subsubsection{Ограничения explainability-слоя}

Knowledge graph и rule-based слой в данной работе играют интерпретационную, а не предиктивную роль. Это важное
достоинство с точки зрения прозрачности архитектуры, но одновременно и ограничение. Explainability-слой:

\begin{itemize}
\tightlist
\item
  не является причинной моделью;
\item
  не доказывает агрономическую обоснованность кандидата;
\item
  не должен трактоваться как независимая экспертная система;
\item
  опирается на компактный набор интерпретационных правил.
\end{itemize}

Следовательно, graph-based explanation не следует понимать как окончательное объяснение «почему именно эта культура
правильна». Он даёт лишь дополнительный семантический контекст для анализа shortlist-кандидатов и позволяет различать
более и менее поддержанные альтернативы. Такой уровень интерпретации полезен для задачи поддержки принятия решений, но
он не заменяет предметную экспертную оценку.

\subsubsection{Ограничения пространственной демонстрации}

Spatial demo и bounding box selection также имеют ограниченный характер. Они не являются отдельной рекомендательной
системой и не создают новую модель; их функция состоит в визуализации и демонстрации уже рассчитанных результатов.
Пространственная карта показывает, как recommendation- и explainability-слои могут быть встроены в пользовательский
интерфейс, однако сама по себе она не расширяет предиктивные возможности системы. Поэтому spatial demo следует
рассматривать как презентационно-аналитический слой, а не как отдельный объект оценки качества.

\subsection{Направления дальнейшего развития}

Несмотря на перечисленные ограничения, полученный прототип задаёт несколько естественных направлений дальнейшего
развития. Эти направления связаны не с полной перестройкой работы, а с логичным расширением уже построенного контура.

\subsubsection{Расширение контекстных признаков}

Одним из наиболее очевидных направлений является добавление новых признаков, отражающих внешний контекст выбора
культуры. В первую очередь к таким признакам можно отнести:

\begin{itemize}
\tightlist
\item
  климатические характеристики;
\item
  почвенные признаки;
\item
  дополнительные региональные агрономические показатели;
\item
  экономические ограничения и предпочтения.
\end{itemize}

В текущей работе подобные источники сознательно не включались в основную постановку, чтобы сохранить воспроизводимость и
управляемую сложность модели. Однако в дальнейшем их использование может усилить прикладную значимость системы и
приблизить её к более реалистичным decision-support сценариям.

\subsubsection{Развитие confidence-анализа}

Ещё одним направлением развития может стать более глубокая работа с оценкой уверенности. В текущей версии
confidence-слой уже показывает, что случаи естественным образом разделяются по уровню надёжности. Однако возможны более
детальные исследования, например:

\begin{itemize}
\tightlist
\item
  уточнение confidence calibration;
\item
  анализ confidence по отдельным классам и регионам;
\item
  построение более формальных uncertainty-aware процедур для shortlist-вывода.
\end{itemize}

Такое развитие позволило бы сделать confidence-компонент не просто полезным дополнением, а более формализованной частью
рекомендательной системы.

\subsubsection{Развитие explainability-слоя}

Knowledge graph в текущей версии уже выполняет роль семантического explainability-слоя, однако его можно развивать в
нескольких направлениях. Во-первых, возможно расширение набора интерпретационных правил. Во-вторых, можно усложнить сам
графовый контекст за счёт введения новых типов сущностей и связей. В-третьих, перспективным представляется сценарий
\textbf{expert-in-the-loop}, в котором пользователь или предметный специалист может корректировать отдельные rule-based
сигналы и анализировать, как это влияет на интерпретацию shortlist-кандидатов.

Такое развитие особенно интересно потому, что оно усиливает именно explainability-компонент работы, не превращая систему
в полноценную экспертную базу знаний. В результате graph-based слой может стать более гибким и более полезным с точки
зрения пользовательского анализа сложных кейсов.

\subsubsection{Развитие пространственного интерфейса}

Пространственная демонстрация уже показывает, как рекомендации могут быть связаны с полями на карте, однако она пока
остаётся исследовательским демо-интерфейсом. В дальнейшем пространственный слой можно развивать как более полноценный
пользовательский инструмент. Возможные направления здесь включают:

\begin{itemize}
\tightlist
\item
  расширение сценариев spatial filtering;
\item
  более удобную навигацию по регионам и полям;
\item
  richer interaction between map, case card and explainability graph;
\item
  интеграцию пространственного интерфейса с rule-based пользовательскими настройками.
\end{itemize}

Такое развитие сделало бы систему ближе к полноценному интерфейсу поддержки принятия решений, сохранив при этом уже
построенное предиктивное и explainability-ядро.

\subsection{Итоговые замечания по ограничениям и развитию}

Перечисленные ограничения не обесценивают полученные результаты, а задают корректные границы их интерпретации.
Разработанная система не претендует на роль полного агрономического оптимизатора, однако в своей целевой постановке она
решает содержательную и практически значимую задачу: формирует shortlist вероятных следующих культур, сопровождает его
confidence-сигналом и дополняет интерпретируемым explainability-контекстом.

Направления дальнейшего развития показывают, что работа может быть расширена как в сторону усиления предиктивного
контекста, так и в сторону развития explainability и интерфейсной части. При этом уже в текущем виде прототип
демонстрирует целостный контур решения и может рассматриваться как основа для дальнейших исследовательских и прикладных
продолжений.
